\documentclass[11pt,a4paper]{article}

\usepackage[utf8]{inputenc}
\usepackage[T1]{fontenc}

\usepackage{amsmath}
\usepackage{amssymb}
\usepackage{amsthm}

\usepackage{graphicx}
\usepackage{geometry}
\usepackage{hyperref}
\usepackage{microtype}

\geometry{margin=2.5cm}

\title{Geometric Reconstruction and Proof of the TIG Horizon Sector}
\author{Kai Stefan Dietrich}
\date{May 11, 2026}

\begin{document}

\maketitle

\begin{abstract}

Topological Integrity Gravity introduced a structurally constrained horizon sector governed by a cubic normal form. In this work, we provide a geometric reconstruction and proof of the TIG1 horizon sector within an effective static and spherically symmetric geometric framework.

Under Schwarzschild asymptotics, regular central structure, one intrinsic correction scale, and admissible horizon transition, the horizon condition reduces to the TIG cubic structure.

The resulting geometry contains:
\begin{itemize}
\item a Schwarzschild-connected outer branch,
\item an inner admissible branch,
\item a degenerate critical horizon,
\item a horizonless supercritical sector.
\end{itemize}

The critical branch merger occurs at

\[
x_c=\frac{2}{3},
\qquad
\beta_c=\left(\frac{4}{27}\right)^{1/3},
\]

and exhibits universal saddle-node scaling behavior.

In addition to the horizon structure, we derive the associated photon-sphere sector and analyze the curvature-invariant structure of the effective geometry.

The present work does not claim complete nonlinear vacuum closure, full tensor consistency, or a final covariant action.

These unresolved problems are deferred to TIG3.

\end{abstract}

\section{Introduction}

Topological Integrity Gravity (TIG) introduced a structurally constrained horizon sector governed by the cubic relation

\[
x^3-x^2+\beta^3=0,
\]

where

\[
x=\frac{r_H}{2M},
\qquad
\beta=\frac{r_c}{2M}.
\]

The original TIG formulation established the existence of a critical horizon structure together with its phenomenological implications. However, the geometric assumptions underlying the cubic sector were not isolated within a fully explicit derivation framework.

The purpose of the present work is to provide a geometric reconstruction and proof of the TIG1 horizon sector under a minimal set of admissibility assumptions.

The analysis is restricted to:
\begin{itemize}
\item static spherical symmetry,
\item Schwarzschild asymptotics,
\item regular central structure,
\item one intrinsic correction scale,
\item admissible horizon transition.
\end{itemize}

Within this effective geometric sector, we show that the horizon condition reduces to the TIG cubic structure.

The resulting geometry contains:
\begin{itemize}
\item a Schwarzschild-connected outer branch,
\item an inner admissible branch,
\item a degenerate critical horizon,
\item a horizonless supercritical sector.
\end{itemize}

The branch merger occurs at the critical point

\[
x_c=\frac{2}{3},
\qquad
\beta_c=\left(\frac{4}{27}\right)^{1/3}.
\]

The resulting structure exhibits universal saddle-node scaling behavior.

In addition to the horizon structure, we derive:
\begin{itemize}
\item the associated photon-sphere sector,
\item the leading TIG correction to the Schwarzschild photon sphere,
\item the curvature-invariant structure of the effective geometry.
\end{itemize}

The resulting effective sector remains curvature-regular while reproducing the Schwarzschild limit asymptotically.

The scope of this paper is deliberately restricted.

The present work does not claim:
\begin{itemize}
\item complete nonlinear vacuum closure,
\item full tensor consistency,
\item a final covariant action,
\item quantization,
\item a complete replacement theory for General Relativity.
\end{itemize}

These unresolved problems are deferred to the next programmatic layer, TIG3.

The goal of TIG2 is therefore the geometric reconstruction and proof of the TIG1 horizon sector within the stated effective geometric framework.

\section{Effective Geometric Sector}

We consider a static and spherically symmetric effective geometry of the form

\[
ds^2
=
-F(r)dt^2
+
\frac{dr^2}{F(r)}
+
r^2d\Omega^2.
\]

The effective mass profile is defined as

\[
m(r)
=
M\frac{r^3}{r^3+r_c^3}.
\]

The corresponding lapse function becomes

\[
F(r)
=
1-\frac{2m(r)}{r}
=
1-\frac{2Mr^2}{r^3+r_c^3}.
\]

This effective sector satisfies:
\begin{itemize}
\item Schwarzschild asymptotics,
\item regular central behavior,
\item one intrinsic correction scale,
\item admissible horizon interpolation.
\end{itemize}

For

\[
r\gg r_c,
\]

the lapse function approaches

\[
F(r)\rightarrow1-\frac{2M}{r},
\]

which reproduces the Schwarzschild limit.

For

\[
r\rightarrow0,
\]

the lapse behaves as

\[
F(r)
\approx
1-\frac{2M}{r_c^3}r^2,
\]

which corresponds to an effective de Sitter-type near-core structure.

The effective geometric sector therefore interpolates continuously between:
\begin{itemize}
\item asymptotic Schwarzschild geometry,
\item and a regularized near-core region.
\end{itemize}

The present work does not claim uniqueness of the effective mass profile.

The role of the effective geometry is restricted to the reconstruction and proof of the TIG horizon sector.

\section{Horizon Reduction}

Horizons are determined by

\[
F(r_H)=0.
\]

Substituting the TIG lapse function gives

\[
1-\frac{2Mr_H^2}{r_H^3+r_c^3}=0.
\]

Rearranging yields

\[
r_H^3+r_c^3-2Mr_H^2=0.
\]

Define the dimensionless variables

\[
x=\frac{r_H}{2M},
\qquad
\beta=\frac{r_c}{2M}.
\]

The horizon equation reduces to

\[
x^3-x^2+\beta^3=0.
\]

This is the TIG horizon cubic.

The cubic therefore follows directly from:
\begin{itemize}
\item static spherical symmetry,
\item Schwarzschild asymptotics,
\item regular central behavior,
\item one intrinsic correction scale,
\item admissible horizon transition.
\end{itemize}

This establishes the geometric reconstruction of the TIG1 horizon sector within the effective geometric framework.

\section{Critical Branch Structure}

Define

\[
P(x,\beta)=x^3-x^2+\beta^3.
\]

The discriminant is

\[
\Delta(\beta)=\beta^3(4-27\beta^3).
\]

The critical point satisfies

\[
P(x_c,\beta_c)=0,
\qquad
\partial_xP(x_c,\beta_c)=0.
\]

This yields

\[
x_c=\frac{2}{3},
\qquad
\beta_c=\left(\frac{4}{27}\right)^{1/3}.
\]

The admissible branch structure is:
\begin{itemize}
\item two admissible positive branches for \(0<\beta<\beta_c\),
\item one degenerate branch at \(\beta=\beta_c\),
\item no admissible horizon branch for \(\beta>\beta_c\).
\end{itemize}

The critical point corresponds to a saddle-node branch merger.

Near criticality, the branch separation obeys

\[
x_\pm-x_c
\sim
(\beta_c-\beta)^{1/2}.
\]

This establishes universal square-root scaling of the TIG horizon sector.

\section{Photon Sphere Structure}

The effective potential for null geodesics is

\[
V_{\mathrm{eff}}(r)
=
\frac{L^2}{r^2}F(r).
\]

Circular null geodesics satisfy

\[
rF'(r)-2F(r)=0.
\]

Substituting the TIG lapse function yields the photon-sphere equation

\[
r^6
-
3Mr^5
+
2r^3r_c^3
+
r_c^6
=
0.
\]

For

\[
r_c\rightarrow0,
\]

the Schwarzschild photon sphere

\[
r_{\mathrm{ph}}=3M
\]

is recovered continuously.

The leading TIG correction shifts the photon sphere inward relative to Schwarzschild.

The photon-sphere structure therefore remains geometrically connected to the TIG horizon sector.

\section{Curvature Invariants}

The effective geometry produces finite curvature invariants throughout the admissible sector.

Near the center,

\[
F(r)
\approx
1-\frac{2M}{r_c^3}r^2,
\]

which corresponds to an effective de Sitter-type structure.

The Ricci scalar remains finite:

\[
R
\rightarrow
\frac{24M}{r_c^3}.
\]

The quadratic Ricci invariant remains finite:

\[
R_{\mu\nu}R^{\mu\nu}
=
\frac{144M^2}{r_c^6}.
\]

The Kretschmann scalar also remains finite:

\[
K
=
\frac{96M^2}{r_c^6}.
\]

For large radius, the geometry approaches the Schwarzschild curvature sector.

The effective geometry therefore preserves:
\begin{itemize}
\item asymptotic Schwarzschild behavior,
\item finite curvature invariants,
\item regular central structure.
\end{itemize}

\section{Structural Meaning}

The TIG horizon sector defines a structurally constrained admissible geometry.

The critical point corresponds to:
\begin{itemize}
\item branch merger,
\item horizon degeneracy,
\item onset of critical scaling,
\item structural transition of the admissible horizon sector.
\end{itemize}

The resulting geometry is not interpreted as a complete gravitational replacement theory.

The present work establishes only the effective geometric reconstruction of the TIG horizon structure.

\section{Claim Boundary}

The present work establishes:
\begin{itemize}
\item geometric reconstruction of the TIG horizon cubic,
\item admissible branch structure,
\item critical horizon merger,
\item photon-sphere structure,
\item curvature regularity of the effective geometric sector.
\end{itemize}

The present work does not establish:
\begin{itemize}
\item complete nonlinear vacuum closure,
\item full tensor consistency,
\item a final covariant action,
\item uniqueness of the effective geometry,
\item quantization,
\item a complete replacement theory for General Relativity.
\end{itemize}

These unresolved problems define the TIG3 program.

\section{Conclusion}

Under static spherical symmetry, Schwarzschild asymptotics, regular central structure, one intrinsic correction scale, and admissible horizon transition, the TIG horizon sector reduces to the cubic relation

\[
x^3-x^2+\beta^3=0.
\]

The resulting geometry contains:
\begin{itemize}
\item admissible horizon branches,
\item a degenerate critical horizon,
\item universal saddle-node scaling,
\item a modified photon-sphere sector,
\item and finite curvature invariants.
\end{itemize}

This establishes the geometric reconstruction and proof of the TIG1 horizon sector within the stated effective geometric framework.

\bibliographystyle{plain}
\bibliography{references}

\end{document}
